


Mumford's deployment of $m$-regularity led to a simplification in the
construction of Quot schemes. The original construction of Grothendieck
had instead relied on Chow coordinates.




\begin{definition}[Castelnuovo--Mumford regularity]\label{nhq-regularity}\source{nitsure-hilbert-quot:page-0009}
Let $k$ be a field and let $\F$ be a coherent sheaf on the
projective space $\PP^n$ over $k$. Let $m$ be an integer.
The sheaf $\F$ is said to be \textbf{$m$-regular} if we have
$$H^i(\PP^n, \F(m-i)) =0\mbox{ for each }i \ge 1.$$
The definition, which may look strange at first sight, is suitable
for making inductive arguments on $n=\dim(\PP^n)$
by restriction to a suitable hyperplane. If $H\subset \PP^n$ is
a hyperplane which does not contain any associated point of $\F$,
then we have a short exact sheaf sequence
$$0 \to \F(m-i-1) \stackrel{\alpha}{\to} \F(m-i) \to \F_H(m-i)\to 0$$
where the map $\alpha$ is locally given by multiplication with
a defining equation of $H$, hence is injective.
The resulting long exact cohomology sequence
$$\ldots \to H^i(\PP^n, \F(m-i)) \to H^i(\PP^n, \F_H(m-i))
\to H^{i+1}(\PP^n, \F(m-i-1)) \to \ldots$$
shows that if $\F$ is $m$-regular, then so is its restriction $\F_H$
(with the \textit{same value\,} for $m$)
to a hyperplane $H\simeq \PP^{n-1}$ which does not contain any
associated point of $\F$.
Note that whenever $\F$ is coherent, the set of associated points of
$\F$ is finite, so there will exist at least one such hyperplane $H$
when the field $k$ is infinite.
\end{definition}

The following lemma is due to Castelnuovo, according to
Mumford's {\sl Curves on a surface}.

\begin{lemma}\label{nhq-castelnuovo}\dcref{2.1}\source{nitsure-hilbert-quot:page-0010}\uses{nhq-regularity}
If $\F$ is an $m$-regular sheaf on $\PP^n$ then
the following statements hold:

\textbf{(a)} The canonical map $H^0(\PP^n,{\mathcal O}_{\PP^n}(1))\otimes
H^0(\PP^n,\F(r)) \to H^0(\PP^n,\F(r+1))$ is surjective
whenever $r\ge m$.

\textbf{(b)} We have $H^i(\PP^n, \F(r))=0$ whenever $i \ge 1$ and $r\ge m-i$.
In other words, if $\F$ is $m$-regular, then it is $m'$-regular for all
$m' \ge m$.


\textbf{(c)} The sheaf $\F(r)$ is generated by its global sections,
and all its higher cohomologies vanish, whenever $r\ge m$.
\end{lemma}

\begin{proof} As the cohomologies base-change correctly under a field extension,
we can assume that the field $k$ is infinite.
We argue by induction on $n$. The statements
\textbf{(a)}, \textbf{(b)} and \textbf{(c)} clearly hold when $n=0$, so
next let $n\ge 1$.
As $k$ is infinite, there exists a hyperplane $H$
which does not contain any associated point of $\F$, so that
the restriction $\F_H$ is again $m$-regular as explained above.
As $H$ is isomorphic to $\PP^{n-1}_k$, by the
inductive hypothesis the assertions of the lemma hold for the sheaf $\F_H$.

When $r=m-i$, the equality $H^i(\PP^n, \F(r))=0$ in statement \textbf{(b)}
follows for all $n\ge 0$ by definition of $m$-regularity.
To prove \textbf{(b)}, we now proceed by induction on $r$ where
$r\ge m-i+1$. Consider the exact sequence
$$H^i(\PP^n, \F(r-1)) \to H^i(\PP^n, \F(r)) \to
H^i(H, \F_H(r))$$
By inductive hypothesis for $r-1$
the first term is zero,
while by inductive hypothesis for $n-1$
the last term is zero, which shows that the middle term is zero,
completing the proof of \textbf{(b)}.

Now consider the commutative diagram
{\footnotesize%scriptsize
$$\begin{array}{lccc}
& H^0(\PP^n, \F(r))\otimes H^0(\PP^n,{\mathcal O}_{\PP^n}(1))&
\stackrel{\sigma}{\to}& H^0(H, \F_H(r))\otimes H^0(H, {\mathcal O}_H(1)) \\
& ~\dna {\scriptstyle \mu}& & ~\dna {\scriptstyle \tau}\\
H^0(\PP^n, \F(r))~~~~ \stackrel{\alpha}{\to} & H^0(\PP^n, \F(r+1))
& \stackrel{\nu_{r+1}}{\to} & H^0(H, \F_H(r+1))
\end{array}$$
}
The top map $\sigma$ is surjective, for the following reason:
By $m$-regularity of $\F$ and using the statement \textbf{(b)}
already proved, we see that
$H^1(\PP^n, \F(r-1))=0$ for $r\ge m$, and so the restriction map
$\nu_r: H^0(\PP^n, \F(r)) \to H^0(H, \F_H(r))$ is surjective.
Also, the restriction map
$\rho : H^0(\PP^n,{\mathcal O}_{\PP^n}(1)) \to H^0(H, {\mathcal O}_H(1))$ is surjective.
Therefore the tensor product $\sigma = \nu_r \otimes \rho$
of these two maps is surjective.

The second vertical map $\tau$ is surjective by inductive
hypothesis for $n-1 = \dim(H)$.

Therefore, the composite $\tau\circ \sigma$ is surjective,
so the composite $\nu_{r+1}\circ\mu$ is surjective,
hence $H^0(\PP^n, \F(r+1)) = \im(\mu) + \ker(\nu_{r+1})$.
As the bottom row is exact, we get
$H^0(\PP^n, \F(r+1)) =  \im(\mu) + \im(\alpha)$.
However, we have $\im(\alpha) \subset \im(\mu)$, as the map
$\alpha$ is given by tensoring with a certain section of
${\mathcal O}_{\PP^n}(1)$ (which has divisor $H$).
Therefore, $H^0(\PP^n, \F(r+1)) = \im(\mu)$.
This completes the proof of \textbf{(a)} for all $n$.

To prove \textbf{(c)}, consider the map $H^0(\PP^n, \F(r))
\otimes H^0(\PP^n,{\mathcal O}_{\PP^n}(p)) \to
H^0(\PP^n, \F(r+p))$,
which is
surjective for $r\ge m$ and $p\ge 0$ as follows from
a repeated use of \textbf{(a)}.
For $p\gg 0$, we know that
$H^0(\PP^n, \F(r+p))$ is generated by its global sections.
It follows that $H^0(\PP^n, \F(r))$ is also generated
by its global sections for $r\ge m$. We already know from \textbf{(b)}
that $H^i(\PP^n, \F(r))=0$ for $i\ge 1$ when $r\ge m$.
This proves \textbf{(c)}, completing the proof of the lemma.
\end{proof}


\begin{remark}\label{nhq-surjectivity-of-restriction}\dcref{2.2}\source{nitsure-hilbert-quot:page-0011}\uses{nhq-castelnuovo}
The following fact, based on the
diagram used in the course of the above proof,
will be useful later: With notation as above,
let the restriction map
$\nu_r: H^0(\PP^n, \F(r)) \to
H^0(H, \F_H(r))$ be surjective. Also, let $\F_H$ be $r$-regular, so that
by Lemma \ref{nhq-castelnuovo}.\textbf{(a)} the map
$H^0(H, {\mathcal O}_H(1))\otimes H^0(H, \F_H(r))\to H^0(H, \F_H(r+1))$
is surjective. Then the restriction map
$\nu_{r+1} : H^0(\PP^n, \F(r+1)) \to  H^0(H, \F_H(r+1))$
is again surjective. As a consequence, if
$\F_H$ is $m$ regular and if for some $r \ge m$ the
restriction map $\nu_r : H^0(\PP^n, \F(r)) \to
H^0(H, \F_H(r))$ is surjective, then the restriction map
$\nu_p : H^0(\PP^n, \F(p)) \to
H^0(H, \F_H(p))$ is surjective for all $p\ge r$.










\end{remark}
\medskip

following theorem.

\begin{theorem}\label{nhq-mumford}\dcref{2.3}\source{nitsure-hilbert-quot:page-0011}\uses{nhq-regularity,nhq-castelnuovo,nhq-surjectivity-of-restriction} {\rm (Mumford)}
For any non-negative integers $p$ and $n$,
there exists a polynomial $F_{p,n}$ in $n+1$ variables
with integral coefficients, which has the following property:

Let $k$ be any field, and let $\PP^n$ denote the $n$-dimensional
projective space over $k$. Let
$\F$ be any coherent sheaf
on $\PP^n$, which is isomorphic to a
subsheaf of $\oplus ^p{\mathcal O}_{\PP^n}$.
Let the Hilbert polynomial of $\F$ be written in terms of binomial
coefficients as
$$\chi(\F(r)) = \sum_{i=0}^n a_i \, {r\choose i}$$
where $a_0,\ldots,a_n \in \Z$.

Then $\F$ is $m$-regular, where %%$m$ is given by the formula
$m=F_{p,n}(a_0,\ldots,a_n)$.
\end{theorem}

\begin{proof} (Following Mumford [M])
As before, we can assume that $k$ is infinite. We argue by induction on $n$.
When $n=0$, clearly we can take $F_{p,0}$ to be any polynomial.
Next, let $n\ge 1$.
Let $H\subset \PP^n$ be a hyperplane which does not contain
any of the finitely many associated points of
$\oplus^p{\mathcal O}_{\PP^n}/\F$ (such an $H$ exists as $k$ is infinite).
Then the following torsion sheaf vanishes:
$$\un{Tor}^1_{{\mathcal O}_{\PP^n}}({\mathcal O}_H,\, \oplus ^p{\mathcal O}_{\PP^n}/\F)=0$$
Therefore the sequence
$0\to \F \to \oplus^p{\mathcal O}_{\PP^n} \to \oplus^p{\mathcal O}_{\PP^n}/\F\to 0$
restricts to $H$ to give a short exact sequence
$0\to \F_H \to \oplus^p{\mathcal O}_H \to \oplus^p{\mathcal O}_H/\F_H \to 0$.
This shows that $\F_H$ is isomorphic to a subsheaf of
$\oplus ^p{\mathcal O}_{\PP^{n-1}_k}$ (under an identification of
$H$ with $\PP^{n-1}_k$),
which is a basic step needed for our inductive argument.

Note that $\F$ is torsion free if non-zero, and so
we have a short exact sequence $0\to \F(-1) \to \F \to \F_H \to 0$.
From the associated cohomology sequence we get
$\chi(\F_H(r)) = \chi(\F(r)) - \chi(\F(r-1))
= \sum_{i=0}^n a_i {r\choose  i} - \sum_{i=0}^n a_i {r-1\choose i}
=  \sum_{i=0}^n a_i {r-1\choose i-1}= \sum_{j=0}^{n-1}b_j{r\choose j}$
where the coefficients $b_0,\ldots,b_{n-1}$ have
expressions $b_j = g_j(a_0,\ldots,a_n)$ where the $g_j$ are polynomials
with integral coefficients independent of the field $k$
and the sheaf $\F$.

By inductive hypothesis on $n-1$
there exists a polynomial $F_{p,n-1}(x_0,\ldots,x_{n-1})$
such that $\F_H$ is $m_0$-regular where
$m_0 = F_{p,n-1}(b_0,\ldots,b_{n-1})$.
Substituting $b_j = g_j(a_0,\ldots,a_n)$, we get
$m_0 = G(a_0,\ldots,a_n)$, where $G$ is a polynomial with
integral coefficients independent of the field $k$
and the sheaf $\F$.

For $m\ge m_0 -1$, we therefore get a long exact cohomology sequence
{\footnotesize
$$0\to H^0(\F(m-1))\to H^0(\F(m)) \stackrel{\nu_m}{\to} H^0(\F_H(m))
\to H^1(\F(m-1))\to H^1(\F(m)) \to 0 \to \ldots $$}
which for $i\ge 2$ gives isomorphisms
$H^i(\F(m-1))\stackrel{\sim}{\to} H^i(\F(m))$. As we have
$H^i(\F(m))=0$ for $m\gg 0$, these equalities show that
$$H^i(\F(m)) = 0\mbox{ for all }i\ge 2\mbox{ and }m\ge m_0-2.$$
The surjections $H^1(\F(m-1))\to H^1(\F(m))$ show that
the function $h^1(\F(m))$
is a monotonically decreasing function of the variable $m$ for $m\ge m_0 - 2$.
We will in fact show that for $m\ge m_0$, the function
$h^1(\F(m))$ is strictly decreasing  till its value reaches zero,
which would imply that
$$H^1(\F(m))=0 \mbox{ for }  m \ge m_0 + h^1(\F(m_0)).$$

Next we will put a suitable upper bound on $h^1(\F(m_0))$ to complete the
proof of the theorem. Note that
$h^1(\F(m-1))\ge h^1(\F(m))$ for $m\ge m_0$,
and moreover equality holds for some $m\ge m_0$
if and only if the restriction map
$\nu_m : H^0(\F(m)) \to H^0(\F_H(m))$ is surjective.
As $\F_H$ is $m$-regular, it follows from
Remark \ref{nhq-surjectivity-of-restriction} that the restriction map
$\nu_j : H^0(\F(j)) \to H^0(\F_H(j))$ is surjective for all
$j\ge m$, so $h^1(\F(j-1)) = h^1(\F(j))$ for all $j\ge m$.
As $h^1(\F(j))=0$ for $j\gg 0$, this establishes our claim that
$h^1(\F(m))$ is strictly decreasing for
$m\ge m_0$ till its value reaches zero.


To put a bound on $h^1(\F(m_0))$, we use the fact that
as $\F \subset \oplus^p{\mathcal O}_{\PP^n}$ we must have
$h^0(\F(r)) \le p h^0({\mathcal O}_{\PP^n}(r)) = p{n + r\choose  n}$.
From the already established fact that
$h^i(\F(m)) = 0$ for all $i\ge 2$ and $m\ge m_0-2$, we now get
\begin{eqnarray*}
h^1(\F(m_0))&=& h^0(\F(m_0)) - \chi(\F(m_0)) \\
            &\le&p{n + m_0\choose  n}-\sum_{i=0}^n a_i {m_0\choose i}\\
            &=& P(a_0,\ldots a_n)
\end{eqnarray*}
where $P(a_0,\ldots, a_n)$ is a polynomial expression in
$a_0,\ldots, a_n$, obtained by substituting
$m_0 = G(a_0,\ldots, a_n)$ in the second line of the above (in)equalities.
Therefore, the coefficients of the corresponding polynomial
$P(x_0,\ldots, x_n)$ are again independent of the field $k$
and the sheaf $\F$. Note moreover that as
$h^1(\F(m_0))\ge 0$, we must have $P(a_0,\ldots, a_n) \ge 0$.


Substituting in an earlier expression, we get
$$H^1(\F(m))=0 \mbox{ for }  m \ge G(a_0,\ldots, a_n) + P(a_0,\ldots, a_n)$$
Taking $F_{p,n}(x_0,\ldots, x_n)$ to be $G(x_0,\ldots, x_n) +
P(x_0,\ldots, x_n)$, and noting the fact that
$P(a_0,\ldots, a_n) \ge 0$, we see that
$\F$ is $F_{p,n}(a_0,\ldots, a_n)$-regular.
This completes the proof of the theorem. \end{proof}

\medskip
